\documentclass[11pt]{article}
\usepackage[margin=2.8cm]{geometry}
\usepackage{newtxtext}
\usepackage[hidelinks]{hyperref}
\pagestyle{empty}
\setlength{\parindent}{0pt}
\setlength{\parskip}{0.6em}

\begin{document}

\begin{flushright}
Amna Dilber\\
Independent researcher, Lahore, Pakistan\\
\texttt{amnadilber.bi@gmail.com}\\
ORCID 0009-0008-5684-4516\\[0.4em]
18 August 2026
\end{flushright}

Dear Editors,

I am submitting ``How much of the evidence for climate-driven dengue
transmission survives the analyst? A multiverse analysis of 221 outbreaks in 33
countries'' for consideration at \emph{Mathematical Biosciences}. I intend to
publish under the subscription model.

Compartmental models fitted to surveillance data are routinely used to test
whether dengue transmission is climate-driven, and reviews explain their
conflicting findings by appealing to local context. Whether that explanation is
right has not been checkable, because each published estimate comes from one
analysis of one place and the two cannot be separated after the fact. This
study separates them: every usable outbreak in a global surveillance
compilation is fitted under all 144 combinations of six analysis choices that
published models rarely state --- 33{,}152 fits.

Roughly three-quarters of the variation in whether climate forcing is endorsed
lies between analyses of the same outbreak rather than between outbreaks.
Partitioned three ways, the largest term is neither: 61\% is an
outbreak-by-analysis interaction against 16\% for the analysis itself. That
distinction decides which remedy is available --- a main effect can be settled
by convention and an interaction cannot --- and it is invisible to a
specification curve, which displays the main effect and averages the
interaction away.

The paper does not stop at the diagnosis. On simulated data whose answer is
known, conventional 95\% intervals contain the truth 71--78\% of the time, and
42\% when the epidemic is a sum of two asynchronous ones, which these data
demonstrably are. Combining eight analyses by Rubin's rules restores nominal
coverage at 1.5--2.2 times the width --- a few minutes of computation, not a
factorial. Applying that rule to the data gives a clear verdict for 137 of 221
outbreaks, 86\% of which support climate forcing: the claim is not that the
effect is absent, but that in 38\% of outbreaks a single season cannot settle
it.

The work fits \emph{Mathematical Biosciences}' stated scope of methodological
articles likely to find application to multiple biological systems. Although the
test case is dengue, the object of study is the estimator and the decision rule:
the decomposition of analytical variation, the failure of nominal interval
coverage under a factorial of defensible model choices, and a validated
multiverse interval apply to any mechanistic model fitted to observational
biological data. The main text is accompanied by a supplement holding the
factor-list robustness analysis, the criterion comparison, further checks and a
worked identifiability case study.

All analyses are reproducible from a public repository containing 38 numbered
pipeline steps, 269 tests --- including automated checks that the manuscript's
quoted figures still match the stored result tables --- and a dated analysis
log recording every formulation that failed and every error found. The
manuscript is available as a preprint on Zenodo
(\url{https://doi.org/10.5281/zenodo.21984527}); it is not under consideration
elsewhere. I have no competing interests and received no funding. I am an
independent researcher without institutional affiliation, and I would be glad
of reviewers willing to be blunt.

Yours sincerely,

Amna Dilber

\end{document}
